\scanpage{225}
\mspage{115}
(15) \unsure{Équivalence|Équivalences} de groupoïdes
\diagram{\begin{tikzcd}[ampersand replacement=\&,column sep=small]
{\begin{gathered}\text{groupoïde des graphes}\\(S,A_0\subset\mathcal P_{\{1,2\}}\deleted{\unsure{6}}(S))\end{gathered}}
\arrow[r,"\approx"] \&
{\begin{gathered}\text{groupoïde des prégraphes}\\\Gamma=(S,S^\vee,R)\text{ séparés en }S^\vee\\
\text{munis d'une jauge }\inserted{\text{surjective}}\ \varphi\end{gathered}}
\end{tikzcd}}
induisant à son tour

(16) Equivalence de groupoïdes \deleted{\unsure{$\longrightarrow$}}
\diagram{\begin{tikzcd}[ampersand replacement=\&,column sep=small]
{\begin{gathered}\text{groupoïde des graphes}\\(S,A_0)\text{ séparés (tout sommet}\\
\text{déterminé par son étoile)}\end{gathered}}\arrow[r,"\approx"] \&
{\begin{gathered}\text{groupoïde des prégraphes}\\\text{biséparés}\inserted{\Gamma},\text{ munis d'}\\
\text{une jauge stricte}\\\text{(i.e. bijective) }\varphi,\\(\text{ou encore (9) d'un}\\
\text{isom }(\varphi,\varphi'):\Gamma\xrightarrow{\sim}\Gamma'\\
\text{tel que }\varphi,\varphi'\text{ inverses l'}\\\text{un de l'autre, où }\Gamma'\\
\text{prégraphe dual) .}\end{gathered}}
\end{tikzcd}}

Soit maintenant
\[
\deleted{\unsure{\lambda}}\deleted{\unsure{\in}} A=\mathrm{Aut}(\Gamma)=
\left\{\inserted{\lambda=}(\lambda_S,\lambda_{S^\vee})\in
\unsure{\mathfrak S|\mathfrak G}_S\times\mathfrak S_{S^\vee}\ \middle|\
\begin{aligned}&\forall(s,s^\vee)\in S\times S^\vee\\
&\deleted{\unsure{\lambda_Ss}}\\
&(s,s^\vee)\in R\Longleftrightarrow\\
&(\lambda_Ss,\lambda_{S^\vee}s^\vee)\in R
\end{aligned}\right\} \tag*{(17)}
\]
\deleted{et \unsure{et} soit $\varphi$ une jauge, d'où une structure de graphe
$\gamma_\varphi=(S,A_\varphi)$. Alors}
\[
\begin{aligned}
&\deleted{\text{Conditions équivalentes}}\\
&\deleted{\text{a) }\lambda_{S^\vee}\varphi=\varphi\lambda_S^{-1}}\\
&\deleted{\text{b) }\lambda_S\text{ est un}}
\end{aligned}\tag*{\deleted{(18)}}
\]
Notons que
\[
\begin{aligned}
&\text{Si }\Gamma\text{ est}\qquad \lambda\longmapsto\lambda_S
  \qquad\text{induit un }\deleted{\text{biject}}\inserted{\text{isom.}}\text{ du groupe}\\
&\text{séparé en }S^\vee\qquad A\overset{\unsure{\supset}}{\longrightarrow}\mathfrak S_S
  \qquad A\text{ avec le groupe des}\\
&\text{i.e. }S^\vee\hookrightarrow\mathcal P(S)
  \hspace{10em}\text{permutations }\deleted{\unsure{\lambda}}f\text{ de }S\\
&\text{alors}\hspace{14em}\text{telles que }\underbrace{f(\mathcal E)}_{\text{tr. de str.}}=\mathcal E\quad\text{i.e.}\\
&\text{où }\mathcal E=\mathrm{Im}(S^\vee\subset\mathcal P(S))\qquad
 \text{qui transforment étoile en}\\
&\hspace{15em}\text{étoile,}
\end{aligned}\tag*{(18)}
\]
\deleted{\unsure{Soit|Soient}} Soit maintenant $\varphi$ une jauge surjective de $\Gamma$ \inserted{séparé en $S^\vee$},
\deleted{définissant donc} \inserted{\unsure{fourn. donc|fournissant|correspondant} \unsure{:|à}} une structure de graphe
\[
 \gamma=\gamma_\varphi=(S,A_\varphi)\quad\text{avec }
 \underbrace{\mathcal E_\gamma}_{\substack{\text{ens des}\\\text{étoiles de }\gamma}}=\mathcal E
 \tag*{(19)}
\]
\deleted{On a alors une} et soit $f\in\mathfrak S_S$. On a alors
% archive sheet number: 224

